\begin{comment}

\section{グループごとに見る自己開示の有効性}
音楽嗜好を通じた自己開示の効果は，以下のような特徴的なグループ間で異なる傾向が観察された．効果が高い～低いまで，3つに分けて詳細を述べる．


\subsubsection{高効果グループ（6ペア，12名）}
% [メモ：I-J, Q-R, S-T, G-H, C-D, K-L]
音楽活動への積極的な参加や日常的な音楽体験を持つグループでは，音楽嗜好の共有を通じた自己開示が特に効果的であった．このグループの被験者は，楽曲に対する具体的な思い入れや経験を豊富に持ち，音楽を通じた対話に対して積極的な態度を示した．特筆すべき点として，楽曲についての技術的な議論から個人的な体験の共有まで，多層的な対話が自然に展開される傾向が見られた．また，相手の知らない楽曲であっても，その楽曲にまつわる個人的な体験や解釈を積極的に共有する様子が観察された．


\subsubsection{中程度効果グループ（3ペア，6名）}
% [メモ：A-B, O-P, L-M]
一般的な音楽愛好者層に相当するグループでは，共通の経験や思い出を持つ楽曲を中心に対話が展開された．このグループの特徴として，段階的に自己開示の深さを増していく傾向が観察された．特に，相手の反応を確認しながら慎重に対話を進める様子が見られ，共通の文脈を持つ楽曲（アニメやドラマの主題歌など）を通じて自己開示が促進される傾向が確認された．

\subsubsection{効果が限定的なグループ（1ペア，2名）}
% [メモ：E-F]
音楽への関与度に大きな差があるペアでは，音楽嗜好の共有を通じた自己開示に困難が見られた．このグループでは，音楽を通じた対話に不安や困難を感じる様子が観察され，楽曲に対する思い入れや経験の共有が限定的となる傾向が確認された．特に，音楽について話すことへの心理的障壁が，自然な対話の展開を妨げる要因となっていた．

\vskip\baselineskip
これらのグループ間での主な差異として，対話の深さ，楽曲選択の特徴，自己開示の質において顕著な違いが観察された．高効果グループでは多層的で深い対話が自然に展開されたのに対し，中程度グループでは段階的な深化が特徴的であり，効果が限定的なグループでは表層的な対話に留まる傾向が見られた．

\end{comment}
